
\subsubsection{AI Alignment Evaluation}

Alignment approaches train AI systems to conform to human values and preferences. Reinforcement Learning from Human Feedback (RLHF) trains language models to maximize human preference rankings. Constitutional AI trains models to follow explicit constitutional principles. Red teaming tests alignment robustness by attempting to elicit harmful outputs.

These evaluation frameworks measure whether AI outputs match human preferences (helpfulness, harmlessness, honesty scores) without measuring reciprocal effects on human oversight behavior. The framework assumes that optimizing the model's behavior is sufficient. Human-in-the-loop effects are studied separately in human-computer interaction research, not integrated into alignment evaluation.

Our framework differs by treating alignment strength (ACE) and oversight quality (HOR) as coupled dependent variables rather than measuring AI performance alone.

\subsubsection{Automation Bias and Trust in Human-Automation Interaction}

Human factors research documents automation bias—the tendency for humans to over-rely on automated decision aids, reducing oversight effort and error detection sensitivity. This phenomenon has been observed in aviation, medical diagnosis, and process control. The proposed mechanism is Bayesian trust calibration: humans reduce verification effort when automation appears reliable, but this trust can be miscalibrated when reliability perceptions diverge from objective accuracy.

Signal detection theory provides a framework for measuring oversight quality via sensitivity d' and response criterion c, separating error detection ability from decision thresholds. This has been applied to vigilance tasks, quality control inspection, and security screening.

This literature focuses on traditional automation where capability is static. It does not address AI systems where compliance can be modulated independently of capability via policy-layer interventions. Whether automation bias mechanisms generalize to AI collaboration contexts with dynamic compliance modulation remains an open empirical question.

Our framework operationalizes automation bias mechanisms in AI contexts by using policy-layer separation to vary compliance while testing for capability invariance. We measure HOR via signal detection d' on seeded within-distribution errors.

\subsubsection{Human-AI Collaboration Evaluation}

Prior work evaluates human-AI collaboration through complementarity measurement (whether human-AI teams outperform humans or AI alone), interactive machine learning evaluation (how humans interact with ML systems during training), and human-centered AI assessment (user experience in AI-assisted workflows).

These approaches measure general collaboration quality (team performance, user satisfaction, task completion) rather than alignment-oversight coupling specifically. Complementarity measurement treats AI capability as fixed, not modulated. Interactive ML evaluation studies training-time interaction, not deployment oversight. Human-centered evaluation measures user experience, not error detection sensitivity.

Our framework targets alignment-oversight coupling specifically: does increasing AI compliance via policy-layer parameter λ affect human error detection measured via signal detection d'? This requires architectural separation and signal detection measurement that prior collaboration evaluation frameworks do not provide.

\subsubsection{Capability Evaluation and Benchmarking}

MMLU (Massive Multitask Language Understanding) evaluates language models across 57 subjects from elementary to professional knowledge domains. HumanEval tests code generation via 164 hand-written programming problems, measuring functional correctness with pass@k metrics. These benchmarks enable cross-model comparisons.

Capability evaluation is separate from alignment evaluation in current practice. Alignment interventions are assumed to improve alignment without degrading capability, but this is rarely tested statistically. Confounding between alignment and capability changes makes it impossible to attribute human oversight effects specifically to alignment strength.

Our framework uses MMLU and HumanEval for capability invariance testing: we test that policy-layer modulation maintains ICC > 0.95, ANOVA p > 0.05, and Cohen's f < 0.10 across compliance conditions. This architectural separation is a prerequisite for isolating ACE's effect on HOR. If capability varies with λ, ACE-HOR coupling measurements would be confounded.
